\section{Discussion}
\label{sec:discussion}

\todo[inline]{Write after the results. Points to cover:
(i) why the anti-alias filter that confines EMI in conventional acquisition is absent by design
below the Nyquist rate;
(ii) what the robust selectors do and do not fix;
(iii) the origin of the Ottawa spikes is unknown, while Paderborn has a known 16~kHz inverter;
(iv) benchmark power as a precondition for claims about learned acquisition;
(v) relation to MRI, where learned masks match equidistant ones at moderate acceleration
\citep{bakker2022adaptive}.}

\paragraph{Limitations}
The study simulates sub-Nyquist acquisition by decimating records acquired at 200 and 64~kS/s. It
models an ideal sampler: no track-and-hold bandwidth, jitter or quantization. The tolerance and the
refinement span are dataset-specific and were fixed before the runs from the geometry and the
record resolution. The Paderborn local ruler exists only where the fault line is visible at full
rate. Energy is reported as converted samples per second, not measured.
